\chapter{Entering Data}
\label{ch:entering}

\section{The two modes}
\index{READY mode}\index{EDIT mode}

\ccprod{} is in one of two states almost all of the time, and the status line
says which.

In \msg{READY} (or \msg{MODIF}, which is \msg{READY} with unsaved changes),
the keys move the cursor and run commands. In \msg{EDIT}, everything you type
goes into a cell.

Typing any ordinary character in \msg{READY} switches to \msg{EDIT} and makes
that character the first one of a new entry. There is no key to press first.

\begin{cccaution}
Typing over a cell replaces it. The old contents are gone the moment you press
\key{Return} --- the first keystroke does not append to what was there. To
change part of a cell rather than replace it, press \key{F2} instead; see
Chapter~\ref{ch:editing}.
\end{cccaution}

\section{Finishing an entry}
\index{committing an entry}

Five keys store what you have typed, and each one also moves the cursor:

\begin{cctable}{Key}{Stores the entry and}{1.55in}{3.27in}
\key{Return} & Moves down one row \\
\key{$\downarrow$} & Moves down one row \\
\key{$\uparrow$} & Moves up one row \\
\key{Tab} & Moves right one column \\
\keyc{Shift}{Tab} & Moves left one column \\
\end{cctable}

\key{Esc} abandons the entry. The cell is left exactly as it was.

\begin{ccnote}
\key{$\leftarrow$} and \key{$\rightarrow$} do \emph{not} finish an entry. While
you are typing they move the insertion point within what you have typed. This
is worth knowing if you are used to a spreadsheet where the left and right
arrows commit and move on.
\end{ccnote}

An entry may be up to eighty characters long. The eighty-first keystroke and
everything after it is discarded without a sound.

\section{What Commander Calc makes of what you typed}
\index{cell!types}

When you finish an entry, \ccprod{} decides what kind of value it is by trying
five things in this exact order.

\begin{ccprocedure}[Order of interpretation]
\item Leading spaces are skipped. If nothing is left, the cell is
\textbf{cleared}.
\item If what remains is \lit{TRUE} or \lit{FALSE}, in any mixture of upper and
lower case, the cell holds a \textbf{boolean}.
\item If the \emph{first} character is \lit{=}, the rest is compiled as a
\textbf{formula}.
\item If the whole entry reads as a number, the cell holds a \textbf{number}.
\item Otherwise the cell holds \textbf{text}.
\end{ccprocedure}

Step~4 is strict: the \emph{entire} entry must be a number. \lit{1.2.3} and
\lit{12A} are not numbers, so they become text. This is what keeps a part
number from turning into arithmetic.

\begin{ccnote}
Step~3 tests the very first character, before the spaces are skipped. So
\lit{=1+1} is a formula and \lit{~=1+1} --- with a leading space --- is a piece
of text that reads \lit{=1+1}. If a cell is showing a formula instead of its
answer, that space is the first thing to look for.
\end{ccnote}

\section{Numbers}
\index{numbers!entering}

A number may have a sign, digits, a decimal point, and an exponent introduced
by \lit{e} or \lit{E}. Spaces before and after are allowed and ignored.

\begin{ccexamples}[These are all numbers]{0.44\linewidth}
\ex{0}{}
\ex{42}{}
\ex{-42}{}
\ex{+42}{}
\ex{123.456}{}
\ex{-0.5}{}
\ex{.25}{}
\ex{8.50}{}
\ex{1e3}{}
\ex{1.5E2}{}
\ex{15e-1}{}
\end{ccexamples}

\begin{ccexamples}[These are text]{0.44\linewidth}
\exn{1.2.3}{}{two decimal points}
\exn{12x}{}{trailing letter}
\exn{1e}{}{no exponent digits}
\exn{-}{}{a sign and nothing else}
\exn{007}{}{a number, in fact --- it becomes 7}
\end{ccexamples}

\begin{ccnote}
The last of those catches people out. \lit{007} \emph{is} a valid number and is
stored as one; the leading zeros are gone the moment you press \key{Return}. If
you need them, type something that is not a number: a leading apostrophe will
not do it, because \ccprod{} has no such convention, but a leading space will.
\end{ccnote}

Numbers are stored to about nine significant digits and displayed, in a cell
with no format, to nine. A number below 1 is shown without a leading zero ---
one third is \lit{.333333333}, in the Commodore tradition --- and a number of
magnitude a thousand million or more is shown in scientific notation.
Appendix~\ref{app:numbers} explains why, and where the limits are.

\section{Text}
\index{text!entering}

Anything the first four steps reject is text. Text may be up to 1,024
characters long, of which the first eighty can be typed --- the rest can only
arrive from a file.

Text is set to the left of its column, and it is cut off at the column edge if
it does not fit. It does not spill over into the empty cell next to it.

\section{TRUE and FALSE}
\index{booleans}

\lit{TRUE} and \lit{FALSE} typed on their own become booleans rather than text,
and they are recognised in any case: \lit{true}, \lit{True} and \lit{TRUE} are
the same thing. They display as \msg{TRUE} and \msg{FALSE} in capitals whatever
you typed.

In arithmetic a boolean counts as 1 or 0, so \fml{=TRUE+TRUE} is 2.

\section{Formulas}
\index{formulas!entering}

An entry beginning with \lit{=} is a formula. \ccprod{} compiles it the moment
you press \key{Return}.

\begin{ccexamples}[Some formulas]{0.44\linewidth}
\ex{=1+1}{2}
\ex{=A1*2}{}
\ex{=SUM(B2:B7)}{}
\ex{=IF(C4>0,"over","under")}{}
\ex{=Q1!B5+Q2!B5}{}
\end{ccexamples}

Chapter~\ref{ch:formularef} is the syntax and Chapter~\ref{ch:funcref} is the functions.

\begin{ccnote}
Only \lit{=} starts a formula. Lotus~1-2-3 users will find that \lit{+A1} and
\lit{@SUM(A1:A9)} are not formulas here; they are text.
\end{ccnote}

\subsection{When a formula will not go in}

If the formula does not compile --- a syntax error, an unknown function name, a
cell reference outside the worksheet, the wrong number of arguments --- then
\ccprod{} \textbf{refuses the entry}. The cell keeps whatever it had before.

\begin{cccaution}
The refusal is silent. No message appears, no error value is stored, and the
cursor moves on as though the entry had worked. What tells you is that the cell
still holds its old value, or is still empty. Check the formula bar.

The reason for the silence is deliberate: the alternative --- storing
\fml{=SUM(A1:A9} as a piece of text because it would not compile --- would turn
a broken formula into a label that looks like it is working.
\end{cccaution}

A formula that compiles but cannot be \emph{evaluated} behaves differently: it
goes into the cell and displays an error value such as \errv{\#DIV/0!}. Those
are listed in Chapter~\ref{ch:messages}.

\section{Dates}
\index{dates!cannot be typed}

\ccprod{} understands dates, stores them, sorts them and displays them as
\lit{YYYY-MM-DD} --- but there is no way to type one. Typing \lit{2026-08-08}
gives you a piece of text; typing \lit{45000} gives you the number 45,000.

Date cells arrive only from an imported \fname{.XLSX} file or a
\fname{.X16S} workbook that already contains them. Chapter~\ref{ch:formats} explains what
happens to them once they are there.

\section{Clearing a cell}
\index{clearing a cell}

Move to it and press \key{Backspace} or \key{Delete}. Either one empties the
cell completely --- value, formula and format. \menupath{Edit > Clear} does the
same.

Pressing \key{Return} on an empty entry clears the cell too, which is the
long way round to the same place.
